\documentclass[11pt]{article}
\usepackage[T1]{fontenc}
\usepackage[utf8]{inputenc}
\usepackage{amsmath,amssymb,amsthm}
\usepackage[margin=2.5cm]{geometry}
\usepackage{booktabs}
\usepackage{graphicx}
\setlength{\parindent}{1.4em}
\setlength{\parskip}{0pt}

\newtheorem{theorem}{Theorem}
\newtheorem{definition}{Definition}
\newtheorem{corollary}{Corollary}
\newtheorem{remark}{Remark}

\begin{document}

% =============== PAGE 1 ===============
\noindent{\LARGE\textbf{Burnout Is Not Too Much Work, v4}}\\[0.3em]
\noindent{\small How to mathematical ragequit your job \hfill 2026}
\vspace{0.3cm}

\noindent{\small\textit{Theoretical model. Calibration illustrative. Action item, not advice.}}
\vspace{0.7cm}

\noindent v1 said burnout is a ratio. v2 said the ratio drifts and the
accounting is observer-dependent. v3 said the denominator decays even at
constant work. All three were diagnosis. v4 is the action item: given
everything we now know, when exactly should you quit?

\medskip\noindent The answer has a name. Charnov 1976 wrote it down for birds
foraging in a patch. The Marginal Value Theorem says a forager should leave
the current patch when its current intake rate equals the long-run average
rate over the cycle of search-plus-forage. The same equation governs jobs.
A bird with no mortgage solves this in seconds. Humans take years.

\begin{center}
\includegraphics[width=0.9\textwidth]{figures/mvt_crossing.pdf}
\end{center}

\medskip\noindent\textbf{Setup.}\quad At time $t$ in a role, person $i$ earns
wage $w$ and pays burnout cost $c(\rho_i(t))$ where $\rho_i(t)$ is the v2+v3
combined ratio. Quitting costs $c_s$ once and forgoes wage for search time
$s$. After quitting, $m_i$ resets to 1 and $\rho_i$ returns to baseline. A
single cycle is one role plus one search. The question is the optimal cycle
length.

\newpage

% =============== PAGE 2 ===============
\begin{definition}[Net Flow Utility]
For person $i$ in role at time $t \geq 0$,
\[
u_i(t) \;=\; w \;-\; c(\rho_i(t)),
\]
where $w > 0$ is wage flow and $c: \mathbb{R}_+ \to \mathbb{R}_+$ is a
continuous, weakly increasing burnout cost. Under the v2+v3 dynamics with
strictly increasing $\rho_i(t)$, $u_i$ is strictly decreasing in $t$.
\end{definition}

\begin{definition}[Long-Run Average Reward]
For a quit time $t^* > 0$ and search time $s > 0$, the long-run average
reward over one role-plus-search cycle is
\[
A_i(t^*) \;=\; \frac{1}{t^* + s} \left[ \int_0^{t^*} u_i(\tau)\, d\tau \;-\; c_s \right],
\]
where $c_s > 0$ is the switch cost and we assume zero income during search.
\end{definition}

\begin{theorem}[Marginal Value Theorem for Job Stopping]
Assume $u_i$ is strictly decreasing and continuous on $[0, \infty)$ with
$u_i(0) > 0$ and $\lim_{t \to \infty} u_i(t) < 0$. Assume $c_s > 0$ and $s > 0$.
Then $A_i$ has a unique interior maximum at $t^* > 0$ characterized by
\[
\boxed{\; u_i(t^*) \;=\; A_i(t^*). \;}
\]
\end{theorem}

\begin{proof}
$A_i$ is differentiable on $(0, \infty)$ with
\[
A_i'(t^*) = \frac{u_i(t^*)(t^* + s) - \left[\int_0^{t^*} u_i\,d\tau - c_s\right]}{(t^*+s)^2}
         = \frac{u_i(t^*) - A_i(t^*)}{t^* + s}.
\]
The sign of $A_i'(t^*)$ matches the sign of $u_i(t^*) - A_i(t^*)$. As
$t^* \to 0^+$, $A_i \to -c_s/s < 0 < u_i(0)$, so $A_i'(0^+) > 0$. As
$t^* \to \infty$, $\int_0^{t^*} u_i\,d\tau \to -\infty$ (because $u_i$ is
eventually negative and decreasing without bound or to a negative
asymptote), so $A_i \to u_i(\infty) < 0$ from above, and $u_i$ continues
decreasing past $A_i$, so eventually $u_i < A_i$ and $A_i' < 0$. By
continuity, $A_i' = 0$ at some $t^*$, where $u_i(t^*) = A_i(t^*)$.
Uniqueness follows because $u_i$ is strictly decreasing while $A_i$ is
single-peaked: once $A_i$ has crossed below $u_i$ from above, it cannot
cross back, since further $t^*$ only adds increasingly negative $u_i$
values to the running average.
\end{proof}

\noindent The proof gives the geometry: $A_i$ rises while it sits below
$u_i$, peaks where they cross, then falls. The peak is the only crossing.

\newpage

% =============== PAGE 3 ===============
\noindent The structurally important fact is that the rational quit time
$t^*$ is \textbf{later} than the threshold crossing $\rho_i(t) = \rho_i^*$,
and \textbf{later} than the moment net utility hits zero. There are three
distinct milestones, and they sit in this order.

\begin{center}
\includegraphics[width=0.85\textwidth]{figures/the_gap.pdf}
\end{center}

\medskip
\begin{tabular}{lll}
\toprule
Milestone & Condition & In our calibration \\
\midrule
Threshold crossing & $\rho_i(t) = \rho_i^*$ & $t = 4.3$ months \\
Subjective hate    & $u_i(t) = 0$ (no net benefit) & $t = 8.5$ months \\
Rational quit      & $u_i(t) = A_i(t)$ (MVT)       & $t^* = 9.1$ months \\
\bottomrule
\end{tabular}

\medskip\noindent The gap between threshold crossing and rational quit (about
$5$ months in our calibration) is not a bug. It is the period during which
amortizing the switch cost still pays off. The subjective sensation is
``something is wrong but I cannot afford to leave yet,'' and that sensation
is exactly correct.

\medskip\noindent\textbf{Worked example.}\quad With $w = 1$, $\rho_0 = 0.6$,
$m_\infty = 0.5$, $\tau_m = 5$ months, environmental drift $\alpha = 0.04$
per month, stress sensitivity $\lambda = 0.7$, switch cost $c_s = 3$ months
of wage, and search time $s = 2$ months, the model gives
\[
t^* \approx 9.1 \text{ months}, \quad \rho_i(t^*) \approx 1.49, \quad u_i(t^*) = A_i(t^*) \approx -0.04.
\]
Note that $A_i(t^*)$ is itself negative. The optimal cycle delivers slightly
negative net welfare. Quitting earlier delivers more negative welfare. Quitting
later delivers more negative welfare. The system is stuck in a regime where
\textit{all} options lose money, and the best move is to lose the least.

\begin{remark}[Why animals do this and humans don't]
The MVT applies whenever a forager has access to an outside option of known
average value. Birds use it because patches are short-lived, search is cheap,
and the brain has dedicated circuitry for tracking running averages. Humans
have long tenure, expensive search, mortgages, employer-tied health insurance,
and no introspective access to $A_i(t^*)$. The math is the same. The
implementation is harder.
\end{remark}

\newpage

% =============== PAGE 4 ===============
\noindent How $t^*$ depends on the parameters. Three things move it.

\begin{center}
\includegraphics[width=0.95\textwidth]{figures/comparative_statics.pdf}
\end{center}

\medskip\noindent Switch cost (left panel) raises $t^*$ monotonically:
expensive transitions force longer cycles. Environmental drift rate $\alpha$
(right panel) lowers $t^*$ monotonically: a deteriorating environment makes
each additional month worse, pulling the optimum forward. Match quality
$m_\infty$ shifts both curves up: better match means longer optimal cycles
because the denominator collapses less.

\medskip\noindent\textbf{The four-lever diagnostic.}\quad If you find yourself
past threshold and contemplating quit, the model says exactly which lever
applies.

\begin{tabular}{ll}
\toprule
What you can change & What it does to $t^*$ \\
\midrule
Reduce switch cost (network, savings, planning) & shifts $t^*$ earlier \\
Reduce search time (better positioned for next role) & shifts $t^*$ earlier \\
Reduce environmental drift (negotiate workload) & shifts $t^*$ later \\
Raise match quality (find a role with higher $m_\infty$) & shifts $t^*$ later or to $\infty$ \\
\bottomrule
\end{tabular}

\medskip\noindent The asymmetry is informative. The first two levers make
quitting easier. The second two make staying viable. There are no levers
that ``do nothing'' to $t^*$. Every action you can take changes when you
should leave.

\medskip\noindent\textbf{Falsification routes.}\quad

\noindent\textit{Route 1.}\quad Predict actual quit timing from individual
$(c_s, s, \alpha, m_\infty)$ estimates. The model says realized quits should
cluster around $t^*$ with deviations explained by liquidity constraints. If
quits are uniform in $t$ or correlate with calendar effects (year-end bonuses),
the model is missing terms.

\noindent\textit{Route 2.}\quad Predict that high-$c_s$ industries (Big Law,
Medicine, Tenure Track, defense contracting) have systematically longer mean
tenures. They do.

\noindent\textit{Route 3.}\quad Predict that workers with shorter expected
search times (high-demand specialties, geographically mobile) quit earlier
controlling for income. Test on labor mobility data.

\newpage

% =============== PAGE 5 ===============
\medskip\noindent\textbf{Empirical anchors.}

\begin{tabular}{lll}
\toprule
Quantity & Value & Source \\
\midrule
Median tenure, US private sector (2024)     & $\sim 3.9$ years     & BLS \\
Mean tenure, Big Law associate              & $\sim 3.0$ years     & NALP \\
Mean tenure, US tech (FAANG)                & $\sim 2.0$ years     & Industry surveys \\
Mean job-search duration (2024)             & $\sim 6$ months      & BLS \\
Switch cost (median total compensation gap) & $\sim 2$--$4$ months & Glassdoor \\
\bottomrule
\end{tabular}

\medskip\noindent The tech tenure number is the diagnostic. Two years is short
even compared to the v3 honeymoon-decay prediction of roughly $1.5$ years to
steady state ($3\tau_m$). Adding the v4 amortization layer pushes optimal
$t^*$ from $1.5$ to roughly $2$ years for typical tech parameters, which
matches the observed mean. Big Law's longer tenure is consistent with high
$c_s$ (specialized human capital, partner-track lock-in). The model matches
the rank ordering of industries.

\medskip\noindent\textbf{The series in one table.}

\begin{tabular}{lll}
\toprule
Post & What it added & What it took out \\
\midrule
v1 & Static threshold $\rho^*$            & n/a \\
v2 & Drift, private threshold, observer wedge & ``rest fixes everything'' \\
v3 & Hedonic decay, denominator dynamics  & ``new job feels different is naivety'' \\
v4 & Optimal stopping, $u(t^*) = A(t^*)$  & ``quit when you cross threshold'' \\
\bottomrule
\end{tabular}

\medskip\noindent\textbf{What this still does not address.}\quad Risk
aversion (the model is risk-neutral; in practice, fear of unemployment biases
$t^*$ later). Job-market state ($A_i(t^*)$ depends on the average value of
\textit{available} next jobs, which fluctuates). Family or care obligations
that lock $c_s$ to a higher value than the worker can change. Anti-pattern
quitting where someone leaves before $t^*$ for emotional reasons and ends up
in a worse role. The model gives the optimum; it does not guarantee
implementation.

\vspace{0.5cm}
\noindent\textbf{References}

\noindent[1] Charnov, E.L. (1976). Optimal foraging, the marginal value theorem. \textit{Theoretical Population Biology} 9(2), 129--136.

\noindent[2] McCall, J.J. (1970). Economics of information and job search. \textit{Quarterly Journal of Economics} 84(1), 113--126.

\noindent[3] Mortensen, D.T. (1986). Job search and labor market analysis. \textit{Handbook of Labor Economics} 2, 849--919.

\noindent[4] Stephens, D.W., Krebs, J.R. (1986). \textit{Foraging Theory}. Princeton University Press.

\noindent[5] Bureau of Labor Statistics (2024). Employee Tenure in 2024. \textit{USDL-24-1881}.

\end{document}
